\chapter{Studiu bibliografic}
\label{cap:studiu_bibliografic}

Acest capitol prezintă stadiul actual al cercetării în domeniile relevante pentru temă: sistemele de recomandare, procesarea imaginilor în domeniul modei și piața soluțiilor existente de asistență vestimentară digitală. Analiza bibliografică ajută la conturarea bazei aplicative pentru proiectul \textit{StyleGenAI}.

\section{Analiza Aplicațiilor Similare Existente}
O cercetare a ofertelor publice pe magazinele de aplicații (Google Play Store și Apple App Store) a relevat existența a două aplicații majore relevante pentru crearea unei idei de pornire și pentru studiul de față: \textbf{Acloset} și \textbf{Whering}.

Ambele aplicații au la bază aceeași idee - digitalizarea garderobei personale pentru a reduce timpul de decizie zilnică și pentru a promova un consum sustenabil. Totuși, abordările și funcționalitățile diferă. Această secțiune își propune să le analizeze obiectiv pentru a stabili limitările ce trebuie depășite.

\subsection{Acloset - AI Fashion Assistant}
Acloset este una dintre cele mai populare aplicații de gestionare a garderobei digitale la nivel global. Creată în Coreea de Sud, aceasta se bazează puternic pe inteligența artificială pentru a automatiza procesul de organizare. Conform datelor publice, aplicația a depășit 1 milion de descărcări.

Aplicația include o funcție de \textit{Auto-Tagging}, unde utilizatorul încarcă o poză, iar sistemul extrage fundalul, recunoscând categoria (ex: tricou, pantaloni) și culoarea. Sistemul sugerează zilnic ținute (OOTD - Outfit of the Day) bazate pe vreme și inventarul disponibil. 

\begin{figure}[H]
    \centering
    \includegraphics[width=0.45\textwidth]{Acloset.jpeg}
    \caption{Interfața aplicației Acloset}
    \label{fig:acloset_app}
\end{figure}

Cu toate acestea, modelul lor s-a orientat masiv către "Platformă de Vânzări Second-Hand". S-a observat empiric că sugestiile generate de AI pe latura stilistică sunt uneori incoerente, combinând articole sport cu elemente formale, demonstrând lipsa unei filtrări semantice riguroase \cite{mcauley2015image}.

\subsection{Whering - Digital Wardrobe \& Styling}
Similar cu Acloset, Whering pune accent pe componenta de sustenabilitate ("Poartă-ți hainele, nu le irosi"). Aplicația permite utilizatorilor să își fotografieze hainele și să planifice ținutele într-un calendar.

\begin{figure}[H]
    \centering
    \includegraphics[width=0.45\textwidth]{w.jpeg}
    \caption{Interfața aplicației Whering}
    \label{fig:whering_app}
\end{figure}

O funcționalitate distinctivă este "Clueless Shuffle", care permite utilizatorului să genereze ținute aleatoare scuturând dispozitivul mobil. Aplicația oferă statistici privind costul amortizat al purtării hainelor (Cost-Per-Wear). Deși interfața este antrenantă, modalitatea sa de potrivire este un proces bazat predominant pe aleatoriu (randomizare) \cite{kang2017visually}. Utilizatorii raportează necesitatea utilizării repetate a funcției shuffle pentru a obține o compatibilitate vizuală acceptabilă.

Concluzionând, deși piața prezintă soluții de top pentru partea de gestiune comercială (Acloset) și sustenabilitate (Whering), abordarea analizei vizuale a calității și compatibilității la \textit{StyleGenAI} vizează un scor algoritmic direct și o armonie cromatică superioară.

\section{Evoluția Asistenților Personali de Stil}
Conceptul de "garderobă digitală" a suferit transformări radicale. Trecerea de la colaje manuale statice la asistenți inteligenți dinamici reflectă maturizarea tehnologiilor de procesare.

\subsection{Era Colajelor Manuale: Polyvore}
Până la mijlocul anilor 2010, platforma \textbf{Polyvore} era responsabilă de generarea inspirației stilistice online. Aceasta permitea utilizatorilor să creeze colaje folosind exclusiv imagini decupate profesional de pre-retaileri. Digitalizarea manuală de către utilizator a propriei garderobe era impracticabilă, implicând utilizarea de editoare grafice externe (ex: Adobe Photoshop) limitând accesul doar utilizatorilor tehnici \cite{vukicevic2018polyvore}.

\subsection{Automatizarea prin AI}
Odată cu implementarea asistentului mobil dotat cu deep learning, bariera extragerii subiectului vizual dispare. Procesul devine un consum fluid în care accentul mută conceptul de la "creație artistică manuală" la "management inteligent". Utilizatorul final abordează un rol pur beneficia din asistența activă (aplicația sugerează proactiv ce să poarte), estompând oboseala decizională (decision fatigue) \cite{fashion_ai_review}.

\section{Stadiul Actual al Tehnologiei: Fashion AI}
Tranziția de la analiza textului asociat hainei la analiza imaginii brute a hainei a cristalizat domeniul "Fashion AI", subdiviziune a Viziunii Artificiale (Computer Vision) \cite{fashion_ai_review}.

Viziunea în domeniu operează de regulă cu:
\begin{enumerate}
    \item \textbf{Analiză Convoluțională}: Rețele tip ResNet50 \cite{he2016deep} extrag ierarhic particularitățile vizuale sub forma seturilor de contururi, texturi și modele compozite.
    \item \textbf{Extracția Semenatică a Subiectelor (Salient Object Detection)}: Modele avansate tip U-Net sau U$^2$-Net \cite{unet_original, qin2020u2net} localizează exact setul de pixeli ce formează obiectul și separă conturul fondului.
    \item \textbf{Căutare după similaritate vizuală}: Construcții neuronale tip \textit{Siamese Networks} transpun imaginea vizuală pe spații paralele decizionale definind matrice de compatibilitate din seturi preexistente \cite{chopra2005learning}.
\end{enumerate}

\section{Sisteme Clasice de Recomandare comparativ cu Fashion}
Soluțiile consacrate utilizează strategii distincte \cite{adomavicius2005toward}:
\begin{itemize}
    \item \textbf{Filtrarea Colaborativă (Collaborative Filtering)}: Asociază istoria altor decizii comune ale masei de conturi \cite{su2009survey}. Totuși, în cazul setărilor garderobelor private, metoda cedează prin fenomenul "Cold Start" (inventarul privat fiind unic).
    \item \textbf{Filtrarea Bazată pe Conținut (Content-Based Filtering)}: Reprezintă motorul fiabil de start, propunând compatibilități axate pe similarități sau reguli de potrivire inter-obiecte pe baza factorilor vizuali (cum procedează și sistemul propus în această lucrare) \cite{lops2011content}.
\end{itemize}

\section{Analiza Comparativă: StyleGenAI vs. Concurenți și Alternative}

Pentru a contextualiza poziția StyleGenAI în piață, se prezintă o comparație structurată cu aplicațiile și metodele existente, pe baza unor criterii relevante.

\subsection{Matrice de Comparație a Aplicațiilor}

\begin{table}[h]
    \centering
    \resizebox{\textwidth}{!}{%
    \begin{tabular}{|l|c|c|c|c|}
    \hline
    \textbf{Criteriu} & \textbf{Acloset} & \textbf{Whering} & \textbf{Polyvore (Legacy)} & \textbf{StyleGenAI (MVP)} \\ \hline
    \multicolumn{5}{|c|}{\textbf{Funcționalități Core}} \\ \hline
    Generare Ținute & Bazată pe AI (black-box) & Aleatoriu ("Shuffle") & Manual (Colaj) & \textbf{Algoritmic (Scorabil)} \\ \hline
    Transparență Algoritm & Scăzută & Scăzută & N/A & \textbf{Ridicată (Funcție de scor)} \\ \hline
    Analiză Culoare & Da (automat) & Da (automat) & Manual & \textbf{Da (K-Means)} \\ \hline
    Integrare Trenduri & Implicită (prin comunitate) & Nu & Manual & \textbf{Da (prin surse de date)} \\ \hline
    \multicolumn{5}{|c|}{\textbf{Arhitectură Tehnică}} \\ \hline
    Backend & Proprietar (Cloud) & Proprietar (Cloud) & Web (Defunct) & \textbf{Python / Flask} \\ \hline
    Frontend & Nativ (iOS/Android) & Nativ (iOS/Android) & Web (Defunct) & \textbf{React Native / Expo} \\ \hline
    Bază de date & Cloud DB & Cloud DB & SQL Server & \textbf{SQLite} \\ \hline
    \multicolumn{5}{|c|}{\textbf{Poziționare și Stadiu}} \\ \hline
    Model Monetizare & Marketplace Second-Hand & Freemium & Publicitate & \textbf{N/A (Propunere: Freemium)} \\ \hline
    Stadiu Proiect & Producție & Producție & Defunct & \textbf{MVP / Prototip Funcțional} \\ \hline
    \end{tabular}%
    }
    \caption{Analiza comparativă a StyleGenAI cu principalii competitori.}
    \label{tab:comparative_analysis}
\end{table}

\subsection{Interpretarea Comparației}

Tabelul \ref{tab:comparative_analysis} evidențiază nișa specifică pe care o ocupă StyleGenAI. În timp ce competitorii se concentrează pe aspecte sociale (Acloset) sau pe gamification (Whering), StyleGenAI se diferențiază prin \textbf{transparența și interpretabilitatea algoritmului}. Utilizatorul nu primește doar o recomandare, ci o recomandare fundamentată pe o funcție de scor explicită, bazată pe teoria culorilor. Această abordare "explainable AI" (XAI) este un avantaj major din punct de vedere academic și pentru utilizatorii care doresc să înțeleagă "de ce"-ul din spatele unei sugestii stilistice.

\subsection{Analiza SWOT pentru StyleGenAI}

O analiză SWOT (Strengths, Weaknesses, Opportunities, Threats) oferă o imagine clară a poziției strategice a proiectului în stadiul său actual.

\begin{table}[h]
    \centering
    \begin{tabular}{|p{0.45\textwidth}|p{0.45\textwidth}|}
    \hline
    \textbf{Puncte Tari (Strengths)} & \textbf{Puncte Slabe (Weaknesses)} \\ \hline
    \begin{itemize}
        \item Algoritm de scoring transparent și explicabil.
        \item Arhitectură modulară (decuplare frontend-backend).
        \item Utilizarea de tehnologii open-source (Python, React Native).
        \item Costuri de rulare zero în stadiul MVP (SQLite).
    \end{itemize} & 
    \begin{itemize}
        \item Scalabilitate limitată de baza de date SQLite.
        \item Lipsa unui mecanism de feedback pentru personalizare (RL).
        \item Stadiu de prototip, fără utilizatori reali pe termen lung.
        \item Clasificare manuală a articolelor, dependentă de utilizator.
    \end{itemize} \\
    \hline
    \textbf{Oportunități (Opportunities)} & \textbf{Amenințări (Threats)} \\ \hline
    \begin{itemize}
        \item Parteneriate de afiliere cu magazine e-commerce.
        \item Extinderea motorului de reguli pentru a include stiluri (formal, casual).
        \item Implementarea personalizării prin Reinforcement Learning.
        \item Integrarea cu tehnologii AR pentru "probă virtuală".
    \end{itemize} &
    \begin{itemize}
        \item Competiție puternică din partea aplicațiilor mature (Acloset, Whering).
        \item Ritmul rapid de schimbare a tendințelor în modă.
        \item Posibila intrare pe piață a giganților tehnologici (Google, Meta).
        \item Dependența de API-uri externe pentru funcționalități viitoare.
    \end{itemize} \\
    \hline
    \end{tabular}
    \caption{Analiza SWOT pentru proiectul StyleGenAI.}
    \label{tab:swot_analysis}
\end{table}

\subsection{Direcții de Dezvoltare Viitoare}
\label{ssec:directii_viitoare}

Analiza comparativă și SWOT subliniază faptul că \textit{StyleGenAI}, în stadiul său de MVP, dispune de o fundație tehnică solidă pe care se pot construi funcționalități avansate. Următoarele direcții reprezintă pași logici în evoluția produsului, transformând slăbiciunile actuale în oportunități strategice:

\begin{enumerate}
    \item \textbf{Scalabilitate la Nivel de Producție:} Cea mai presantă nevoie tehnică este migrarea de la SQLite la un sistem de baze de date robust, capabil să gestioneze un număr mare de utilizatori și operațiuni concurente. Soluții precum \textbf{PostgreSQL} sau un serviciu cloud (ex: Amazon RDS, Google Cloud SQL) ar oferi scalabilitatea, securitatea și performanța necesare pentru o lansare publică.

    \item \textbf{Personalizare prin Reinforcement Learning (RL):} Pentru a depăși motorul de reguli static, se poate implementa un model de RL. Sistemul ar putea învăța preferințele stilistice unice ale fiecărui utilizator analizând acțiunile acestuia: ținutele salvate, cele șterse sau cele modificate manual. Astfel, recomandările ar deveni cu adevărat personalizate și s-ar adapta în timp.

    \item \textbf{Integrare E-commerce și Monetizare:} Oportunitatea majoră de business constă în crearea de parteneriate cu retaileri online. Prin integrarea API-urilor acestora, aplicația ar putea sugera utilizatorilor articole vestimentare noi, complementare garderobei existente. Acest model de afiliere (commission-based) reprezintă o cale directă de monetizare.
    
    \item \textbf{Clasificare Automată a Articolelor:} Pentru a reduce efortul utilizatorului, se poate dezvolta un model de clasificare a imaginilor (bazat pe o rețea neuronală convoluțională, ex: MobileNetV2) care să identifice automat categoria articolului (tricou, pantaloni, etc.) după încărcarea fotografiei.
\end{enumerate}

Aceste dezvoltări ar poziționa \textit{StyleGenAI} nu doar ca un instrument utilitar, ci ca o platformă completă de management și inspirație vestimentară.

\section{Concluzia Studiului Bibliografic}

Literele și analiza pieței confirmă că domeniul fashion tech este dinamic și competitiv. StyleGenAI contribuie prin:
\begin{enumerate}
    \item Claritate și interpretabilitate algoritm (vs. black-box competitors).
    \item Integrare trend semantică automată.
    \item Performance acceptabilă pe MVP-uri mobile.
    \item Fundație stabilă pentru extensii viitoare (AR, RL, e-commerce).
\end{enumerate}

Succesul sa va depinde de capacitate de scaling, parteneri e-commerce și evoluție Features în linie cu preferințele pieței.